% ============================== ANNEXE ============================
\settab{LECTURE VI}{8}
\chapter[Multiple Integrals (Lecture VI)]{Multiple Integrals\\ (Lecture VI)}\label{ch:six}\chaptermark{Multiple Integrals}

\begin{examnote}
Lecture VI (2 November) has no TD sheet and we do not know whether it is examined. If it is,
expect computations rather than theory: an iterated integral, a change of order, a polar,
cylindrical or spherical integral, an area, a volume or a centre of mass. The theory of Darboux
sums is summarised but the problems focus on calculation.
\end{examnote}

\section{Definition and properties}

\begin{intuition}
Cut the region into small rectangles, multiply the height $f$ by the area of each, and add: a
Riemann sum. As the rectangles shrink the sum tends to the volume under the graph (or the mass,
if $f$ is a density). Lower and upper sums trap the true value from below and above.
\end{intuition}

\begin{defi}[Double integral]
For bounded $f$ on a rectangle $R$ and a partition $P$, $L(P,f)=\sum_S(\inf_Sf)\,\mathrm{Area}(S)$ and
$U(P,f)=\sum_S(\sup_Sf)\,\mathrm{Area}(S)$. $f$ is integrable if $\sup_PL(P,f)=\inf_PU(P,f)$; the
common value is $\iint_Rf\dd A$. Over a general region $A\subset R$, integrate the extension of $f$
by $0$ outside $A$.
\end{defi}

\begin{thm}[Darboux criterion and integrable functions]
$f$ is integrable iff for every $\varepsilon>0$ some partition has $U(P,f)-L(P,f)<\varepsilon$.
Continuous functions on a closed rectangle are integrable, and so are bounded functions that
are continuous except on finitely many smooth curves. The integral is linear, monotone and
additive over regions that overlap only along their boundaries.
\end{thm}

\begin{trap}
Bounded does not mean integrable: the function equal to $1$ on $\mathbb Q^2$ and $0$ elsewhere has
$L=0$ and $U=1$ for every partition.
\end{trap}

\section{Fubini and the order of integration}

\begin{thm}[Fubini on a rectangle, regions between graphs]
$\iint_{[a,b]\times[c,d]}f=\int_a^b\!\int_c^df(x,y)\dd y\dd x=\int_c^d\!\int_a^bf(x,y)\dd x\dd y$ for continuous $f$.
On $A=\{a\le x\le b,\ g_1(x)\le y\le g_2(x)\}$: $\iint_Af=\int_a^b\int_{g_1(x)}^{g_2(x)}f\dd y\dd x$.
\end{thm}

\begin{method}
\textbf{Setting up limits.} (1) Sketch the region. (2) Choose the order: vertical slices
($dy$ inside) or horizontal slices ($dx$ inside). (3) For the inner integral, read where a slice
enters and leaves the region; the outer limits are constants. (4) If the inner antiderivative
is impossible (like $e^{y^2}$), describe the same region in the other order.
\end{method}

\begin{center}
\begin{tikzpicture}
\begin{axis}[bp, width=5.2cm, height=5.2cm, xmin=-0.15, xmax=1.25, ymin=-0.15, ymax=1.25,
  xtick={1}, ytick={1}, xlabel=$x$, ylabel=$y$, axis equal image]
  \addplot[name path=lo, tealD, domain=0:1, samples=40] {x^2};
  \addplot[name path=hi, coralD, domain=0:1] {x};
  \addplot[aqua, fill opacity=0.35] fill between[of=lo and hi];
  \draw[tealDD, line width=1pt, -{Stealth}] (axis cs:0.6,0.36) -- (axis cs:0.6,0.6);
\end{axis}
\end{tikzpicture}\hspace{5mm}
\begin{tikzpicture}
\begin{axis}[bp, width=5.2cm, height=5.2cm, xmin=-0.15, xmax=1.25, ymin=-0.15, ymax=1.25,
  xtick={1}, ytick={1}, xlabel=$x$, ylabel=$y$, axis equal image]
  \addplot[name path=d, tealD, domain=0:1] {x};
  \addplot[name path=t, draw=none, domain=0:1] {1};
  \addplot[blush, fill opacity=0.5] fill between[of=d and t];
  \draw[coralD, line width=1pt, -{Stealth}] (axis cs:0,0.7) -- (axis cs:0.7,0.7);
\end{axis}
\end{tikzpicture}
\figcap{Left: the region between $y=x^2$ and $y=x$, sliced vertically. Right: $0\le x\le y\le1$, the
region of $\int_0^1\int_x^1e^{y^2}\dd y\dd x$, sliced horizontally after the change of order.}
\end{center}

\section{Polar, cylindrical and spherical coordinates}

\begin{thm}[Change of variables, informal general form]
For a smooth one-to-one map $G$ with nonzero Jacobian inside $B$:
$\iint_{G(B)}f=\iint_Bf(G(u))\,|\det DG(u)|\dd u$.
\begin{center}
\begin{tabular}{@{}lll@{}}
polar & $(r\cos\theta,r\sin\theta)$ & $dA=r\dd r\dd\theta$\\
cylindrical & $(r\cos\theta,r\sin\theta,z)$ & $dV=r\dd z\dd r\dd\theta$\\
spherical & $(\rho\sin\varphi\cos\theta,\rho\sin\varphi\sin\theta,\rho\cos\varphi)$ & $dV=\rho^2\sin\varphi\dd\rho\dd\varphi\dd\theta$
\end{tabular}
\end{center}
with $\varphi\in[0,\pi]$ measured from the positive $z$-axis.
\end{thm}

\begin{remember}
Useful polar curves: $r=a\cos\theta$ ($-\frac\pi2\le\theta\le\frac\pi2$) is the circle of radius $\frac a2$
centred at $(\frac a2,0)$; $r=a\sin\theta$ ($0\le\theta\le\pi$) the one centred at $(0,\frac a2)$. Centre of
mass: $\bar x=\frac1m\iint xf\dd A$ with $m=\iint f\dd A$.
\end{remember}

\begin{reallife}
The Gaussian integral: $I=\int_{-\infty}^{\infty}e^{-x^2}dx$ satisfies
$I^2=\iint_{\R^2}e^{-(x^2+y^2)}dA=\int_0^{2\pi}\!\!\int_0^\infty e^{-r^2}r\dd r\dd\theta=\pi$, so $I=\sqrt\pi$: the
normalisation of the normal distribution in probability.
\end{reallife}

\section{Exam-style problems for Lecture VI}

\begin{problem}{Lecture VI · iterated integrals}
(a) $\int_1^2\!\int_{-3}^4x^2y\dd y\dd x$. (b) $\iint_A(x^2+y^2)\dd A$, $A$ between $y=x^2$ and $y=x$.
(c) The area of $A$.
\end{problem}
\begin{solution}
(a) Inner: $x^2\big[\tfrac{y^2}2\big]_{-3}^4=\tfrac72x^2$; outer: $\tfrac72\cdot\tfrac73=\tfrac{49}6$.
(b) $\int_0^1\big[x^2y+\tfrac{y^3}3\big]_{x^2}^{x}\dd x=\int_0^1\big(x^3+\tfrac{x^3}3-x^4-\tfrac{x^6}3\big)\dd x
=\tfrac14+\tfrac1{12}-\tfrac15-\tfrac1{21}=\tfrac3{35}$.
(c) $\int_0^1(x-x^2)\dd x=\tfrac16$.
\end{solution}

\begin{problem}{Lecture VI · change of order}
Compute $I=\int_0^1\!\int_x^1e^{y^2}\dd y\dd x$.
\end{problem}
\begin{hint}$e^{y^2}$ has no elementary antiderivative in $y$. Redescribe $0\le x\le y\le1$.\end{hint}
\begin{solution}
The region is $0\le y\le1$, $0\le x\le y$, so $I=\int_0^1\!\int_0^ye^{y^2}\dd x\dd y=\int_0^1ye^{y^2}\dd y
=\tfrac12\big[e^{y^2}\big]_0^1=\tfrac{e-1}2$.
\end{solution}

\begin{problem}{Lecture VI · splitting a region}
$\iint_A2xy\dd A$ over the triangle bounded by $y=0$, $y=x$, $x+y=2$.
\end{problem}
\begin{solution}
Split at $x=1$: $\int_0^1\int_0^x2xy\dd y\dd x=\int_0^1x^3\dd x=\tfrac14$ and
$\int_1^2\int_0^{2-x}2xy\dd y\dd x=\int_1^2x(2-x)^2\dd x=\tfrac5{12}$. Total $\tfrac23$. (In the other order,
$0\le y\le1$, $y\le x\le2-y$, one integral suffices.)
\end{solution}

\begin{problem}{Lecture VI · polar}
(a) $\iint_D(x^2+y^2)\dd A$ over the unit disc. (b) $\iint_{x^2+y^2\le4}e^{-(x^2+y^2)}\dd A$.
(c) $\iint x^2\dd A$ inside the cardioid $r=1-\cos\theta$.
\end{problem}
\begin{solution}
(a) $\int_0^{2\pi}\!\int_0^1r^2\cdot r\dd r\dd\theta=2\pi\cdot\tfrac14=\tfrac\pi2$.
(b) $2\pi\int_0^2e^{-r^2}r\dd r=\pi(1-e^{-4})$.
(c) $\int_0^{2\pi}\!\int_0^{1-\cos\theta}r^3\cos^2\theta\dd r\dd\theta=\tfrac14\int_0^{2\pi}(1-\cos\theta)^4\cos^2\theta\dd\theta=\tfrac{49\pi}{32}$.
\end{solution}

\begin{problem}{Lecture VI · spherical and centre of mass}
(a) Volume inside $x^2+y^2+z^2=1$ and above the cone $z^2=x^2+y^2$, $z\ge0$.
(b) Centre of mass of the uniform quarter disc $x,y\ge0$, $x^2+y^2\le1$.
(c) $\bar z$ for the uniform first-octant part of the unit ball.
\end{problem}
\begin{solution}
(a) The cone is $\varphi=\frac\pi4$: $V=\int_0^{2\pi}\!\int_0^{\pi/4}\!\int_0^1\rho^2\sin\varphi\dd\rho\dd\varphi\dd\theta
=2\pi\cdot\tfrac13\big(1-\tfrac{\sqrt2}2\big)=\tfrac{\pi(2-\sqrt2)}3$.
(b) $m=\frac\pi4$, $\iint x\dd A=\int_0^{\pi/2}\!\int_0^1r^2\cos\theta\dd r\dd\theta=\tfrac13$, so
$\bar x=\bar y=\frac{4}{3\pi}$.
(c) $\iiint z\dd V=\int_0^{\pi/2}\!\int_0^{\pi/2}\!\int_0^1\rho^3\cos\varphi\sin\varphi\dd\rho\dd\varphi\dd\theta=\frac\pi{16}$,
$m=\frac\pi6$, so $\bar z=\frac38$.
\end{solution}

\begin{problem}{Lecture VI · triple integral}
For the tetrahedron $T$: $x,y,z\ge0$, $x+y+z\le1$, compute $\mathrm{Vol}(T)$ and $\iiint_Tx\dd V$.
\end{problem}
\begin{solution}
$\int_0^1\!\int_0^{1-x}\!\int_0^{1-x-y}\dd z\dd y\dd x=\tfrac16$; with $x$ inside, $\tfrac1{24}$. So the centre of
mass has $\bar x=\frac{1/24}{1/6}=\frac14$ (and by symmetry $\bar y=\bar z=\frac14$).
\end{solution}

\begin{trap}
The most common integration errors: forgetting $r$ in $dA=r\dd r\dd\theta$ (or $\rho^2\sin\varphi$); outer
limits that depend on the inner variable; and splitting a region but integrating one piece twice.
\end{trap}

% ------------------------------------------------------------------

\section{Quick check}
\begin{quickcheck}
For continuous $f$ on a rectangle, the two iterated integrals agree. & \qT & Fubini.\\
In polar coordinates $dA=dr\,d\theta$. & \qF & $dA=r\dd r\dd\theta$.\\
In spherical coordinates $dV=\rho^2\sin\varphi\dd\rho\dd\varphi\dd\theta$. & \qT & Determinant of the spherical map.\\
$\int_0^1\!\int_0^xf\dd y\dd x=\int_0^1\!\int_0^yf\dd x\dd y$. & \qF & Same region is $y\le x\le1$: $\int_0^1\!\int_y^1f\dd x\dd y$.\\
A bounded function continuous except on a curve of area $0$ is integrable. & \qT & Lecture VI criterion.\\
$\int_{-\infty}^{\infty}e^{-x^2}\dd x=\sqrt\pi$. & \qT & Square it and go polar.\\
\end{quickcheck}

\settab{LECTURE VII}{9}
\chapter[Change of Variables (Lecture VII)]{The Change of Variables\\ Formula (Lecture VII)}\label{ch:seven}\chaptermark{Change of Variables}

\section{Determinants measure area and volume}

\begin{thm}[Area and volume]
The parallelogram spanned by $A=(a,c)$, $B=(b,d)$ has area $|ad-bc|$. The box spanned by
$A,B,C$ in $\R^3$ has volume $|\det(A,B,C)|$. A linear map $L$ multiplies every area (volume)
by $|\det L|$.
\end{thm}

\begin{intuition}
Area $=|A|\,|B|\,|\sin\theta|$, and $|A|^2|B|^2-(A\cdot B)^2=(ad-bc)^2$. For a general region,
approximate it by small squares: each one's area is multiplied by $|\det L|$, and the boundary
squares have total area going to $0$.
\end{intuition}

\begin{center}
\begin{tikzpicture}[scale=0.9]
  \draw[bpgrid, step=0.5] (-0.5,-1.5) grid (6,2.5);
  \draw[tealDD, ->] (-0.5,0) -- (6,0); \draw[tealDD, ->] (0,-1.5) -- (0,2.5);
  \fill[aqua, opacity=0.25] (0,0) rectangle (1,1); \draw[tealD] (0,0) rectangle (1,1);
  \node[font=\scriptsize, text=tealD] at (0.5,0.5) {$S$};
  \fill[coral, opacity=0.3] (0,0) -- (2,1) -- (5,0) -- (3,-1) -- cycle;
  \draw[coralD, line width=1pt] (0,0) -- (2,1) -- (5,0) -- (3,-1) -- cycle;
  \draw[vec] (0,0) -- (2,1) node[above, font=\scriptsize]{$(2,1)$};
  \draw[vec] (0,0) -- (3,-1) node[below, font=\scriptsize]{$(3,-1)$};
  \node[font=\small, text=coralD] at (2.6,0) {$L(S)$};
\end{tikzpicture}
\figcap{The unit square maps to the parallelogram spanned by $(2,1)$ and $(3,-1)$: area
$|2\cdot(-1)-1\cdot3|=5$.}
\end{center}

\begin{remember}
Dilations: $(x,y)\mapsto(rx,sy)$ multiplies areas by $rs$; so the ellipse $\frac{x^2}{a^2}+\frac{y^2}{b^2}\le1$
has area $\pi ab$ and the ellipsoid with semi-axes $a,b,c$ has volume $\frac43\pi abc$. A tetrahedron
spanned by $A_1,A_2,A_3$ has volume $\frac16|\det(A_1,A_2,A_3)|$; with four arbitrary vertices,
translate one to the origin first.
\end{remember}

\section{The formula}

\begin{thm}[Change of variables]
If $G$ is $C^1$-invertible on the interior of $R$ and $f$ is continuous on $G(R)$ except on finitely
many smooth curves (surfaces), then
\[ \iint_{G(R)}f(x,y)\dd x\dd y=\iint_Rf(G(u,v))\,|\Delta_G(u,v)|\dd u\dd v,\qquad
\Delta_G=\det\begin{pmatrix}\partial_ug_1&\partial_vg_1\\ \partial_ug_2&\partial_vg_2\end{pmatrix}, \]
and the same in $\R^3$ with $3\times3$ determinants.
\end{thm}

\begin{method}
\textbf{Choosing a substitution.} (1) Pick $u,v$ that make the region a rectangle or the
integrand simple (e.g.\ $u=x+y$, $v=x-y$). (2) Express $x,y$ in terms of $u,v$. (3) Compute
$\Delta_G=\det\partial(x,y)/\partial(u,v)$ and take its absolute value. (4) Transform the region,
then integrate.
\end{method}

\section{Exam-style problems for Lecture VII}

\begin{problem}{Lecture VII · determinants}
(a) Area of the parallelogram spanned by $(3,1)$ and $(2,-5)$. (b) Area of the triangle with
vertices $(1,2)$, $(3,-1)$, $(0,0)$. (c) Volume of the parallelotope spanned by $(3,1,2)$, $(1,-1,4)$,
$(2,1,0)$. (d) Volume of the tetrahedron spanned by $(3,1,4)$, $(-1,2,1)$, $(5,-2,1)$.
\end{problem}
\begin{solution}
(a) $|3\cdot(-5)-1\cdot2|=17$. (b) $\tfrac12|1\cdot(-1)-3\cdot2|=\tfrac72$.
(c) $\det=3(0-4)-1(0-8)+2(1+2)=2$, volume $2$ (the slide's $-4$ is a slip).
(d) $\det=3(2+2)-1(-1-5)+4(2-10)=-14$, volume $\tfrac{14}6=\tfrac73$.
\end{solution}

\begin{problem}{Lecture VII · ellipse and ellipsoid}
Areas and volumes of $\frac{x^2}9+\frac{y^2}{16}\le1$ and $\frac{x^2}9+\frac{y^2}{16}+\frac{z^2}{25}\le1$. Image of the
ellipse $\frac{x^2}{a^2}+\frac{y^2}{b^2}\le1$ under $L=\begin{pmatrix}1&2\\3&-5\end{pmatrix}$.
\end{problem}
\begin{solution}
$(u,v)\mapsto(3u,4v)$ maps the unit disc onto the ellipse: area $12\pi$. $F_{3,4,5}$ maps the unit ball onto
the ellipsoid: volume $3\cdot4\cdot5\cdot\frac43\pi=80\pi$. $|\det L|=|-5-6|=11$, so the image has area $11\pi ab$.
\end{solution}

\begin{problem}{Lecture VII · nonlinear map}
$G(u,v)=(e^u\cos v,e^u\sin v)$ on $0\le u\le1$, $0\le v\le\pi$. Compute $\iint_{G(R)}x^2\dd A$.
\end{problem}
\begin{solution}
$\Delta_G=e^{2u}$ and $x^2=e^{2u}\cos^2v$, so the integral is
$\int_0^1e^{4u}\dd u\int_0^\pi\cos^2v\dd v=\frac{e^4-1}4\cdot\frac\pi2=\frac{\pi(e^4-1)}8$. ($G(R)$ is the upper half
of the annulus $1\le r\le e$.)
\end{solution}

\begin{problem}{Lecture VII · a rotated square}
Compute $\iint_S(x+y)^2\dd A$ over $S=\{|x|+|y|\le1\}$ with $u=x+y$, $v=x-y$.
\end{problem}
\begin{solution}
$S$ becomes the square $-1\le u,v\le1$. Since $x=\frac{u+v}2$, $y=\frac{u-v}2$,
$\Delta=\det\begin{pmatrix}\frac12&\frac12\\ \frac12&-\frac12\end{pmatrix}=-\frac12$, so $|\Delta|=\frac12$ and
$\iint_Su^2\cdot\frac12\dd u\dd v=\frac12\cdot\frac23\cdot2=\frac23$.
\end{solution}

\begin{trap}
Use $|\Delta_G|$, not $\Delta_G$: orientation does not matter for area. And when the new variables
are given as $u=u(x,y)$, either invert them or use $\Delta_{(x,y)/(u,v)}=1/\Delta_{(u,v)/(x,y)}$.
\end{trap}

\section{Quick check}
\begin{quickcheck}
The parallelogram spanned by $A,B$ has area $\det(A,B)$. & \qF & Area is $|\det(A,B)|$.\\
A linear map $L$ multiplies every area by $|\det L|$. & \qT & Theorem B.1.\\
The tetrahedron spanned by $A,B,C$ has volume $|\det(A,B,C)|$. & \qF & One sixth of it.\\
The change of variables formula uses $|\Delta_G|$ evaluated on $R$, the source region. & \qT & $\iint_{G(R)}f=\iint_R(f\circ G)|\Delta_G|$.\\
The region $\frac{x^2}{a^2}+\frac{y^2}{b^2}\le1$ has area $\pi ab$. & \qT & $G(u,v)=(au,bv)$, $\det=ab$.\\
A negative Jacobian makes the integral negative. & \qF & The formula uses the absolute value.\\
\end{quickcheck}
